\chapter{Sources de lumière et propagation}\label{ch:g7:light-sources-propagation}

Pourquoi peux-tu lire cette page~? Une lampe brille --- mais la page
n'est pas une lampe, et pourtant ton œil reçoit de la \omterm{def:g1:light-and-shadows:source}{lumière} de
chaque ligne imprimée. Entre les objets qui fabriquent de la \omterm{def:g1:light-and-shadows:source}{lumière}
et ceux qui se contentent de la transmettre se trouve une distinction
dont l'optique ne peut pas se passer~; et avec elle vient la réponse
complète à la plus vieille question de ce livre~: que faut-il pour
voir quelque chose~?

\section{Fabricants et transmetteurs}

\begin{definition}[Sources primaires]\label{def:g7:light-sources-propagation:primary}
Une \emph{source primaire}\index{source primaire} fabrique sa propre
\omterm{def:g1:light-and-shadows:source}{lumière}~: le Soleil et les \omterm{def:g4:solar-system:star}{étoiles}, les flammes, les filaments
incandescents, l'acier chauffé à blanc de la forge, les lucioles, les
éclairs. Sa \omterm{def:g1:light-and-shadows:source}{lumière} existe que quoi que ce soit d'autre brille ou
non --- éteins le monde, une source primaire brille encore.
\end{definition}

\begin{definition}[Objets diffusants]\label{def:g7:light-sources-propagation:diffusion}
Tout autre objet visible est un \emph{objet
diffusant}\index{diffusion}~: il reçoit de la \omterm{def:g1:light-and-shadows:source}{lumière} et en renvoie
--- il en \emph{diffuse} --- une partie dans toutes les directions.
La page, la pomme, la \omterm{def:g2:moon-in-the-sky:moon}{Lune}, ton propre visage~: aucun ne fabrique la
moindre lueur~; chacun disperse généreusement de tous côtés la
\omterm{def:g1:light-and-shadows:source}{lumière} qu'il reçoit. Un objet diffusant, dans le noir, est tout
simplement invisible.
\end{definition}

\begin{omfigure}
\begin{tikzpicture}[scale=0.8]
  % the lamp, a primary source
  \fill[omThm] (0.7,2.7) circle (0.3);
  \foreach \a in {0,45,...,315} {
    \draw[omThm, thick] (0.7,2.7) ++(\a:0.42) -- ++(\a:0.2);
  }
  \node[above, font=\small] at (1.1,3.35) {lampe~: source primaire};
  % its light travels to the apple
  \draw[-Stealth, omThm, thick] (1.35,2.75) -- (3.55,2.15);
  \draw[-Stealth, omThm, thick] (1.35,2.45) -- (3.45,1.65);
  \draw[-Stealth, omThm, thick] (1.25,2.15) -- (3.35,1.15);
  % the apple scatters it to every side
  \draw[thick, fill=omMeth, fill opacity=0.5] (5.2,1.2) circle (0.5);
  \foreach \a in {15,60,105,150,195,240,285,330} {
    \draw[-Stealth, omDef, thick] (5.2,1.2) ++({\a}:0.6)
      -- ++({\a}:1.05);
  }
  \node[below, font=\small] at (5.2,-0.5) {pomme~: objet diffusant};
  \node[right, font=\small, omDef] at (7.35,2.15)
    {la lumière diffusée part de tous côtés~:};
  \node[right, font=\small, omDef] at (7.35,1.55)
    {chaque œil de la pièce en reçoit};
\end{tikzpicture}

{\small La division du travail~: la lampe fabrique la \omterm{def:g1:light-and-shadows:source}{lumière}~; la
pomme, qui la reçoit, en disperse une part dans toutes les
directions --- et voilà pourquoi tout le monde, autour de la table,
voit la même pomme.}
\end{omfigure}

\begin{example}[Diffuser n'est pas
réfléchir]\label{ex:g7:light-sources-propagation:notmirror}
Les \omterm{def:g5:mirrors-reflection:reflection}{miroirs} comme les pages renvoient la \omterm{def:g1:light-and-shadows:source}{lumière} reçue --- mais avec
des manières opposées. Le \omterm{def:g5:mirrors-reflection:reflection}{miroir} la fait rebondir dans \emph{une}
direction disciplinée (angles égaux --- ta vieille règle du rebond)
en préservant l'\omterm{prop:g5:mirrors-reflection:image}{image}~; la rugosité microscopique de la page projette
la \omterm{def:g1:light-and-shadows:source}{lumière} \emph{partout}, ce qui détruit l'\omterm{prop:g5:mirrors-reflection:image}{image} mais l'offre à tous
les yeux à la fois. Voilà pourquoi un \omterm{def:g5:mirrors-reflection:reflection}{miroir} montre la lampe, alors
que la page ne montre qu'elle-même, d'où que tu la regardes.
\end{example}

\begin{example}[Le dossier de la Lune, clos]\label{ex:g7:light-sources-propagation:moon}
Le verdict de l'enfance a maintenant ses mots officiels~: la \omterm{def:g2:moon-in-the-sky:moon}{Lune}
n'est pas une \omterm{def:g7:light-sources-propagation:primary}{source primaire}~; c'est un \omterm{def:g7:light-sources-propagation:diffusion}{objet diffusant} de taille
planétaire --- de la roche grise qui disperse la \omterm{def:g1:light-and-shadows:source}{lumière} du Soleil
dans toutes les directions, la nôtre comprise. Le clair de lune est
donc du soleil de seconde main, et la «~face sombre~» est simplement
la moitié dont le dossier ne contient, à cet instant, aucune \omterm{def:g1:light-and-shadows:source}{lumière}
à transmettre.
\end{example}

\section{La vision, réglée}

\begin{proposition}[Les deux conditions de la
vision]\label{prop:g7:light-sources-propagation:vision}
Un œil voit un objet exactement quand
\begin{enumerate}
  \item de la \omterm{def:g1:light-and-shadows:source}{lumière} \emph{quitte} l'objet --- fabriquée, si
        l'objet est une \omterm{def:g7:light-sources-propagation:primary}{source primaire}~; diffusée, s'il est un
        transmetteur --- et que
  \item cette \omterm{def:g1:light-and-shadows:source}{lumière} voyage en ligne droite, sans obstacle, jusque
        dans l'œil.
\end{enumerate}
Casse l'une ou l'autre condition --- un \omterm{def:g7:light-sources-propagation:diffusion}{objet diffusant} non éclairé,
ou un mur en travers de la ligne de vue --- et l'objet n'est pas vu.
Toutes les énigmes de visibilité de ce livre se ramènent à ces deux
contrôles.
\end{proposition}

\begin{example}[Le faisceau rendu visible]\label{ex:g7:light-sources-propagation:beam}
Dans un \omterm{def:g2:air-around-us:air}{air} propre, un \omterm{def:g6:light-propagation:ray}{faisceau} de lampe de poche qui traverse la
pièce ne se voit pas de côté --- l'\omterm{def:g2:air-around-us:air}{air} ne transmet presque rien, et
aucune \omterm{def:g1:light-and-shadows:source}{lumière} du \omterm{def:g6:light-propagation:ray}{faisceau} ne se tourne vers ton œil. Secoue un
chiffon poussiéreux~: le \omterm{def:g6:light-propagation:ray}{faisceau} surgit. Chaque grain de poussière
est un minuscule diffuseur qui attrape la \omterm{def:g1:light-and-shadows:source}{lumière} du \omterm{def:g6:light-propagation:ray}{faisceau} et t'en
disperse une part --- des milliers de grains dessinant pour toi la
ligne droite du \omterm{def:g6:light-propagation:ray}{faisceau}. Le brouillard, la brume et la fumée rendent
le même service aux faisceaux des phares et aux projecteurs de
cinéma.
\end{example}

\section{Trois habits pour la lumière}

\begin{definition}[Transparent, translucide,
opaque]\label{def:g7:light-sources-propagation:coats}
Les matériaux rencontrent la \omterm{def:g1:light-and-shadows:source}{lumière} de trois manières~:
\begin{enumerate}
  \item \emph{transparent}\index{transparent}~: la \omterm{def:g1:light-and-shadows:source}{lumière} traverse
        presque sans être dérangée et garde son ordre --- verre
        clair, eau immobile, \omterm{def:g2:air-around-us:air}{air}~; les \omterm{prop:g5:mirrors-reflection:image}{images} passent intactes~;
  \item \emph{translucide}\index{translucide}~: la \omterm{def:g1:light-and-shadows:source}{lumière}
        traverse, mais brouillée --- verre dépoli, papier calque,
        rideau léger~; la clarté passe, les \omterm{prop:g5:mirrors-reflection:image}{images} non~;
  \item \emph{opaque}\index{opaque}~: la \omterm{def:g1:light-and-shadows:source}{lumière} est arrêtée ---
        absorbée ou diffusée en arrière --- bois, pierre, métal, ta
        main~; derrière un objet opaque s'étend une \omterm{def:g1:light-and-shadows:shadow}{ombre}.
\end{enumerate}
\end{definition}

\begin{method}[Trier des matériaux à la lampe de
poche]\label{met:g7:light-sources-propagation:sort}
Une pièce sombre, une lampe de poche, une page imprimée, et
l'échantillon entre les deux~:
\begin{enumerate}
  \item tenir l'échantillon au-dessus du texte, la lampe derrière
        lui~;
  \item on lit le texte à travers~? --- \omterm{def:g7:light-sources-propagation:coats}{transparent}~;
  \item on voit une lueur mais aucune lettre~? --- \omterm{def:g7:light-sources-propagation:coats}{translucide}~;
  \item on voit du noir, et une \omterm{def:g1:light-and-shadows:shadow}{ombre} au-delà~? --- \omterm{def:g7:light-sources-propagation:coats}{opaque}.
\end{enumerate}
Fais l'essai honnêtement~: bien des matériaux changent d'équipe avec
l'épaisseur --- une couche d'eau est transparente, la mer profonde
est noire comme la nuit~; une feuille de papier est \omterm{def:g7:light-sources-propagation:coats}{translucide}, le
livre fermé est \omterm{def:g7:light-sources-propagation:coats}{opaque}.
\end{method}

\begin{omfigure}
\begin{tikzpicture}[scale=0.8]
  \foreach \x/\name in {0.5/transparent, 4.4/translucide,
      8.3/opaque} {
    \fill[omThm] (\x,1.5) circle (0.22);
    \draw[thick] (\x+1.3,0.4) rectangle (\x+1.7,2.6);
    \node[below, font=\small] at (\x+1.6,-0.0) {\name};
  }
  \draw[-Stealth, omThm, thick] (0.75,1.5) -- (1.78,1.5);
  \draw[-Stealth, omThm, thick] (2.22,1.5) -- (3.4,1.5);
  \draw[-Stealth, omThm, thick] (4.65,1.5) -- (5.68,1.5);
  \foreach \a in {28,10,-12,-30} {
    \draw[-Stealth, omThm, thin] (6.12,1.5) -- ++({\a}:1.1);
  }
  \draw[-Stealth, omThm, thick] (8.55,1.5) -- (9.58,1.5);
  \draw[line width=3pt, omExo] (10.02,0.55) -- (10.02,2.45);
  \node[right, font=\footnotesize, omExo] at (10.1,0.8) {côté ombre};
\end{tikzpicture}

{\small Trois sorts pour un rayon~: passer en bon ordre~; passer mais
brouillé~; être arrêté, avec une \omterm{def:g1:light-and-shadows:shadow}{ombre} derrière.}
\end{omfigure}

\begin{example}[Les trois équipes de
l'architecture]\label{ex:g7:light-sources-propagation:architecture}
Les bâtisseurs déploient les trois à dessein. Le verre \omterm{def:g7:light-sources-propagation:coats}{transparent}
des fenêtres~: la vue et la \omterm{def:g1:light-and-shadows:source}{lumière} ensemble. Les vitres \omterm{def:g7:light-sources-propagation:coats}{translucides}
des salles de bain et les cloisons dépolies des bureaux~: la \omterm{def:g1:light-and-shadows:source}{lumière}
du jour entre, l'intimité est préservée --- la clarté sans les
\omterm{prop:g5:mirrors-reflection:image}{images}, telle est exactement leur fiche de poste. Les murs et volets
\omterm{def:g7:light-sources-propagation:coats}{opaques}~: l'obscurité à la demande. Les abat-jour sont de la
translucidité au service du confort~: l'éblouissement du filament
incandescent y est brouillé en une lueur douce et sans source
apparente.
\end{example}

\begin{remark}[Rien à transmettre, rien à
voir]\label{rem:g7:light-sources-propagation:space}
Devinette~: une astronaute entre les \omterm{def:g4:solar-system:star}{étoiles} regarde \emph{de côté}
un vide traversé de rayons solaires. Que voit-elle~? Du noir --- alors
que la \omterm{def:g1:light-and-shadows:source}{lumière} du Soleil inonde ses yeux en passant. L'espace ne
porte pas de poussière digne de ce nom~: rien ne diffuse vers elle la
\omterm{def:g1:light-and-shadows:source}{lumière} qui passe, et une \omterm{def:g1:light-and-shadows:source}{lumière} qui n'entre pas dans l'œil ne peint
rien. Le ciel du jour est bleu précisément parce que notre \omterm{def:g2:air-around-us:air}{air}, lui,
transmet vers nous la \omterm{def:g1:light-and-shadows:source}{lumière} du Soleil --- toute une \omterm{def:g4:solar-system:planet}{planète} de
douce \omterm{def:g7:light-sources-propagation:diffusion}{diffusion} au-dessus de nos têtes. Pas d'\omterm{def:g2:air-around-us:air}{air}, ciel noir, même à
midi~: les astronautes de la \omterm{def:g2:moon-in-the-sky:moon}{Lune} se tenaient en plein soleil sous
une coupole noire et sans \omterm{def:g4:solar-system:star}{étoiles}.
\end{remark}

\section{Exercices}

\begin{exercise}[$\star$]\label{exo:g7:light-sources-propagation:1}
Trie~: la flamme d'une bougie~; \ la \omterm{def:g2:moon-in-the-sky:moon}{Lune}~; \ une luciole~; \ cette
page~; \ un feu rouge~; \ la \omterm{def:g4:solar-system:planet}{planète} Vénus au crépuscule.
\end{exercise}

\begin{exercise}[$\star$]\label{exo:g7:light-sources-propagation:2}
Que fait un \omterm{def:g7:light-sources-propagation:diffusion}{objet diffusant} de la \omterm{def:g1:light-and-shadows:source}{lumière} qu'il reçoit --- et que
devient-il dans une pièce plongée dans le noir~?
\end{exercise}

\begin{exercise}[$\star$]\label{exo:g7:light-sources-propagation:3}
Énonce les deux conditions de la vision, et sers-t'en~: pourquoi ne
peut-on pas voir la surface \emph{elle-même} d'un \omterm{def:g5:mirrors-reflection:reflection}{miroir}
parfaitement propre, mais seulement les objets qu'il réfléchit~?
(Attention --- belle question~; réponds simplement~: que refuse de
faire son poli avec la \omterm{def:g1:light-and-shadows:source}{lumière}~?)
\end{exercise}

\begin{exercise}[$\star$]\label{exo:g7:light-sources-propagation:4}
Pourquoi tout le monde, autour de la table, voit-il la même pomme,
alors qu'une seule position de l'œil, chanceuse, attrape le reflet de
la lampe dans une cuillère~?
\end{exercise}

\begin{exercise}[$\star$]\label{exo:g7:light-sources-propagation:5}
Pourquoi un \omterm{def:g6:light-propagation:ray}{faisceau} de lampe de poche est-il invisible de côté dans
un \omterm{def:g2:air-around-us:air}{air} propre, et visible dans la brume~? Nomme le phénomène.
\end{exercise}

\begin{exercise}[$\star$]\label{exo:g7:light-sources-propagation:6}
Classe~: le verre à vitre~; \ le papier calque~; \ une brique~; \ une
faible épaisseur d'eau claire~; \ l'enveloppe d'une ampoule dépolie~;
\ du papier d'aluminium.
\end{exercise}

\begin{exercise}[$\star$]\label{exo:g7:light-sources-propagation:7}
Quel métier fait la translucidité dans une fenêtre de salle de bain
et dans un abat-jour~? Qu'est-ce qui passe, qu'est-ce qui est
arrêté~?
\end{exercise}

\begin{exercise}[$\star$]\label{exo:g7:light-sources-propagation:8}
La \omterm{def:g2:moon-in-the-sky:moon}{Lune} est éclairée par le Soleil --- et pourtant les astronautes
lunaires voyaient un ciel noir à midi. Explique avec
\cref{rem:g7:light-sources-propagation:space}, et dis pourquoi le
ciel de midi de la Terre est bleu vif à la place.
\end{exercise}

\begin{exercise}[$\star\star$]\label{exo:g7:light-sources-propagation:9}
Une seule feuille de papier luit chaudement devant une lampe de
poche~; le livre fermé arrête entièrement la \omterm{def:g1:light-and-shadows:source}{lumière}. Que fait
l'épaisseur à l'appartenance d'un matériau à une équipe~? Donne un
second exemple de matériau qui change d'équipe.
\end{exercise}

\begin{exercise}[$\star\star$]\label{exo:g7:light-sources-propagation:10}
Astuce de cinéma~: les faisceaux des projecteurs sont invisibles dans
une salle propre, et pourtant les films les montrent tranchant
dramatiquement l'\omterm{def:g2:air-around-us:air}{air}. Que l'équipe de tournage doit-elle ajouter à
l'\omterm{def:g2:air-around-us:air}{air}, et quelle physique achète-t-elle ainsi~?
\end{exercise}

\begin{exercise}[$\star\star$]\label{exo:g7:light-sources-propagation:11}
Tu vois le visage d'un camarade près d'un feu de camp. Retrace le
voyage complet de la \omterm{def:g1:light-and-shadows:source}{lumière} --- source, \omterm{def:g7:light-sources-propagation:diffusion}{diffusion}, lignes droites
--- et compte combien de diffuseurs se trouvent dans la chaîne si tu
vois \emph{en outre} ce visage réfléchi dans le lac.
\end{exercise}

\begin{exercise}[$\star\star\star$]\label{exo:g7:light-sources-propagation:12}
Le spectacle lumineux de la nature~: par temps clair, le ciel loin du
Soleil est bleu, le disque du Soleil d'un blanc aveuglant --- et au
coucher, le Soleil bas devient orange tandis que les nuages
au-dessus se teintent de rose. En n'utilisant que «~l'\omterm{def:g2:air-around-us:air}{air} diffuse un
peu de la \omterm{def:g1:light-and-shadows:source}{lumière} qui le traverse, et d'autant plus que le trajet est
long~», explique pourquoi le soleil bas est dépouillé d'une partie de
sa \omterm{def:g1:light-and-shadows:source}{lumière} --- et qui reçoit la part dérobée. (L'histoire complète
des couleurs sera le chapitre sur les spectres de l'an prochain~;
raisonne ici uniquement avec plus-de-trajet, plus-de-diffusion.)
\end{exercise}

\section{Problème~: éclairer le musée}

\begin{problem}[{Problème du week-end --- la nouvelle salle des
sculptures~; la nuit de travail d'une conceptrice d'éclairage~;
sources, diffuseurs et trois sortes de
verre}]\label{pb:g7:light-sources-propagation:1}
La nouvelle salle du musée ouvre lundi~: une statue de marbre, des
tableaux aux murs, une vitrine d'exposition, et un conservateur
nerveux. Tu assistes la conceptrice d'éclairage.

\medskip\noindent\textbf{Partie I --- les principes, sur place.}
\begin{enumerate}
  \item Le conservateur demande pourquoi la statue, «~si
        éclatante de blancheur~», a malgré tout besoin de lampes.
        Réponds avec les deux conditions de la vision.
  \item Énumère les \omterm{def:g7:light-sources-propagation:primary}{sources primaires} et les \omterm{def:g7:light-sources-propagation:diffusion}{objets diffusants} de
        la salle une fois les lampes allumées~: les lampes, la
        statue, les tableaux, les visiteurs, le verre de la
        vitrine.
  \item Le marbre semble luire de partout, sans aucune tache
        éblouissante. Quelle manière de renvoyer la \omterm{def:g1:light-and-shadows:source}{lumière} est à
        l'œuvre, et pourquoi toutes les positions de visiteur en
        reçoivent-elles~?
  \item Une plaque de laiton poli, en revanche, renvoie un unique
        reflet éblouissant d'une lampe. Nomme l'autre manière, et
        énonce la règle qui gouverne l'endroit où part cet éclat.
\end{enumerate}

\medskip\noindent\textbf{Partie II --- apprivoiser l'éblouissement.}
\begin{enumerate}[resume]
  \item Les ampoules nues éblouissent les visiteurs. La conceptrice
        enferme chacune dans un globe dépoli. À quelle équipe de
        matériaux appartient le globe, et que fait-il à la \omterm{def:g1:light-and-shadows:source}{lumière}
        du filament~?
  \item Les tableaux doivent être éclairés, mais aucune lampe ne
        doit apparaître \emph{réfléchie} dans leur vernis depuis
        l'endroit d'où on les regarde. À l'aide de la règle du
        \omterm{def:g5:mirrors-reflection:reflection}{miroir}, où la conceptrice place-t-elle les lampes par
        rapport au tableau et au spectateur~?
  \item Le verre clair de la vitrine est presque invisible ---
        jusqu'à ce que la manche blanche d'un visiteur y apparaisse
        comme un fantôme. Pourquoi un verre clair transmet-il
        \emph{et} réfléchit-il faiblement~? Qu'obtient la
        conceptrice avec sa niche aux parois noires derrière la
        vitrine~?
  \item Des grains de poussière dérivent parfois dans les faisceaux
        des projecteurs, dessinant des fils lumineux dans l'\omterm{def:g2:air-around-us:air}{air}. Le
        conservateur~: «~Magnifique~! Gardez-le.~» La conceptrice~:
        «~Cela signifie que l'\omterm{def:g2:air-around-us:air}{air} est sale.~» Qui a raison sur la
        physique, et que sont ces fils~?
\end{enumerate}

\medskip\noindent\textbf{Partie III --- l'essai de nuit.}
\begin{enumerate}[resume]
  \item \omterm{def:g1:light-and-shadows:source}{Lumières} éteintes, volets scellés~: le gardien rapporte une
        salle «~absolument noire --- même pas la statue blanche~».
        Explique avec les deux conditions pourquoi la blancheur est
        impuissante dans le noir.
  \item Un puits de \omterm{def:g1:light-and-shadows:source}{lumière} est proposé~: du verre dépoli au
        plafond. Que livrera-t-il le jour, et que refusera-t-il de
        livrer (la vue des nuages, l'éblouissement direct du
        soleil)~?
  \item Le conservateur craint que le soleil ne fasse pâlir les
        tableaux et demande «~de la \omterm{def:g1:light-and-shadows:source}{lumière} sans aucun soleil~». La
        conceptrice sourit~: de jour, même la douce \omterm{def:g1:light-and-shadows:source}{lumière} des
        fenêtres orientées au nord \emph{est} du soleil.
        Explique --- qui l'a transmise~?
  \item Rédige le rapport de la conceptrice~: trois phrases ---
        une sur la raison pour laquelle la \omterm{def:g7:light-sources-propagation:diffusion}{diffusion} est la
        meilleure amie du musée, une sur les endroits où il faut
        surveiller les reflets à la manière du \omterm{def:g5:mirrors-reflection:reflection}{miroir}, et une sur
        l'équipe de verre choisie pour chaque service.
\end{enumerate}
\end{problem}
